\documentclass[../main.tex]{subfiles}

\begin{document}
    \subsection{Nama Program}
    \begin{lstlisting}[language=C++]
        #include <iostream>
    \end{lstlisting}

    {\Script} ``{\crnfamily \#include <iostram>}'' adalah \textit{header file} yang digunakan sebagai standar masukan dan keluaran dalam bahasa C++ untuk mendeklarasikan dan mendefinisikan {\statement} “{\crnfamily cin}” dan “{\crnfamily cout}”.

    \begin{lstlisting}[language=C++]
        using namespace std;
    \end{lstlisting}

    {\Script} ''{\crnfamily using namespace std;}'' merupakan {\statement} yang digunakan untuk memberi tahu {\compiler} bahwa seluruh variabel dan fungsi yang berada dalam {\scope} {\namespace} ``{\crnfamily std}'' akan menjadi bagian dari {\scope} file saat ini menghilangi kebutuhan untuk mengetik ``{\crnfamily std::}'' sebelum memanggil variabel atau fungsi dalam {\namespace} ``{\crnfamily std}''.
\end{document}
